\section{Theoretical results}

\subsection{An exact two-crop reversal frontier}

Consider two crops with \(s_1>s_2\), fixed planted land
\(x_1+x_2=A\), and an operationally feasible interval
\(I=[\underline z,\,\overline z]\) for \(z=x_1\). Under a declared
joint-law path \(\theta\), write
\[
r_\theta(z)=\CVaR_\alpha[-z\widetilde\pi_1-(A-z)\widetilde\pi_2],
\qquad
\mathcal K_{\theta,\kappa}=\{z\in I:r_\theta(z)\le\kappa\}.
\]

\begin{theorem}[Necessary-and-sufficient restricted reversal]
\label{thm:frontier}
Suppose \(\mu_1>\mu_2\) and \(\mathcal K_{\theta,\kappa}\) is nonempty.
Then the unique optimum is
\[
z^*(\theta,\kappa)=\sup\mathcal K_{\theta,\kappa}.
\]
Consequently,
\[
x_1^*<x_2^*
\Longleftrightarrow
\sup\mathcal K_{\theta,\kappa}<A/2
\Longleftrightarrow
\mathcal K_{\theta,\kappa}\cap[A/2,\,\overline z]=\varnothing .
\]
If \(r_\theta\) is nondecreasing on
\([A/2,\,\overline z]\) and a risk-feasible point exists below \(A/2\),
reversal is equivalent to \(r_\theta(A/2)>\kappa\).
\end{theorem}

\paragraph{Proof.}
The objective is
\(\mu_2A+(\mu_1-\mu_2)z\), which is strictly increasing in \(z\).
Loss-CVaR is convex and continuous in \(z\), so its sublevel set intersected
with compact \(I\) is compact. The strictly increasing objective is uniquely
maximized at the right endpoint. The three reversal statements then follow
from \(x_1<x_2\Longleftrightarrow z<A/2\). Under upper-half monotonicity,
the lowest risk on that half is attained at \(A/2\), giving the last
equivalence. \(\square\)

The theorem applies to a fixed-land two-crop model with ordered expected
margins, operational bounds and a nonempty risk-feasible set. If
\(r_\theta(A/2)\) is continuous and strictly increasing along a specified
dependence path and crosses \(\kappa\), there is a unique family-specific
threshold \(\theta^*\) solving
\(r_{\theta^*}(A/2)=\kappa\). Without strict monotonicity or crossing, the
reversal set is
\(\{\theta:\sup\mathcal K_{\theta,\kappa}<A/2\}\); it may contain several
components. Higher tail dependence therefore need not produce a monotone or
universal reversal response.

\subsection{Full-model optimality and selection}

\begin{theorem}[Complete convex KKT certificate]
\label{thm:kkt}
Under a convex constraint qualification, or for the bounded finite-scenario
linear programme, a feasible solution of Eq.~\ref{eq:canonical} is optimal if
and only if there exist nonnegative multipliers for every inequality,
together with the free-threshold and scenario-excess dual variables, that
satisfy primal feasibility, dual feasibility, stationarity
Eq.~\ref{eq:kkt} and complementarity.
\end{theorem}

The proof follows convex KKT sufficiency or finite-dimensional linear
programming strong duality. Incomplete stationarity can assign a residual
margin difference to ``risk'' even when budget, contract or shared capacity is
active. KKT provides pointwise optimality conditions; optimal-face analysis
determines whether the ranking relation holds throughout all equivalent
optima.

\subsection{Dependence paths}

\begin{proposition}[Named-family feasible-set ordering]
\label{prop:convexorder}
Fix all crop marginals. If
\(L_{\theta_1}(x)\preceq_{\mathrm{cx}}L_{\theta_2}(x)\) for every
\(x\) in a declared domain, then
\[
\CVaR_{\alpha,\theta_1}[L(x)]
\le \CVaR_{\alpha,\theta_2}[L(x)].
\]
If the domain is all of \(\mathcal X\), the risk-feasible set under
\(\theta_2\) is contained in that under \(\theta_1\), so optimal expected
profit cannot increase. Individual crop shares need not be monotone.
\end{proposition}

Loss-CVaR is monotone in convex order, yielding the risk and feasible-set
inclusions. But set inclusion does not order coordinates of different
optimizers. We therefore report Gaussian, Student-\(t\) and Clayton results
family by family, including active-set transitions, multiple optima and
disconnected regions.

\subsection{An executable diversification-failure criterion}

Let \(x^{MV}\) denote the \emph{selected Gaussian mean--variance policy}: the
first point on the specified Gaussian frontier that attains a 15\%
variance-reduction target relative to the expected-profit benchmark \(x^0\).
Let \(x^T\) denote the tail-aware policy under the Student-\(t\) evaluation
law.

\begin{definition}[Diversification failure]
Conventional diversification fails when all three conditions hold:
\begin{align}
\operatorname{Var}_G[\Pi(x^{MV})]&<
\operatorname{Var}_G[\Pi(x^0)],\\
\CVaR_{\alpha,T}[L(x^{MV})]&>
\CVaR_{\alpha,T}[L(x^T)],\\
x^{MV}&\ne x^T.
\end{align}
Here, \(x^0\) is the expected-profit policy under the matched Gaussian law.
In computation, allocation disagreement requires
\(\lVert x^{MV}-x^T\rVert_1>\varepsilon_x\), and tail inferiority requires
\(\CVaR_{\alpha,T}[L(x^{MV})]>
\CVaR_{\alpha,T}[L(x^T)]+\varepsilon_r\).
Failure is operationally strong when
\(\CVaR_{\alpha,T}[L(x^{MV})]>\kappa+\varepsilon_r\).
\end{definition}

\begin{proposition}
If \(x^T\) is feasible and the preceding three conditions hold,
\(x^{MV}\) is not tail-risk optimal. Under the strong inequality it also
violates the declared risk ceiling.
\end{proposition}

The criterion separates Gaussian variance reduction, allocation disagreement
and downside-risk infeasibility under the heavy-tail evaluation law. It
applies to the specified benchmark comparison rather than every diversified
portfolio. If
\(\CVaR_{\alpha,T}[L(x^T)]=\kappa\), the tail-inferiority and ceiling-violation
inequalities coincide numerically and must not be counted as independent
evidence.

\subsection{Information, complementarity and substitution}

A state \(\omega\) is realized and produces a noisy signal \(S\), while a
fraction of the planting plan remains adjustable. The producer observes
\(S\), chooses a signal-contingent action and subsequently realizes margin.
Loss-CVaR is evaluated ex ante over the combined state--signal distribution.

\begin{theorem}[Restricted information value]
In a posterior-separable model, or when the shared ex-ante risk restriction is
nonbinding, if ignoring the signal is admissible, optimized information value
is nonnegative. It is zero when one common action is optimal for every
posterior. If posterior optima are unique and differ on positive-probability
signal events, with strict improvement on at least one event, information
value is strictly positive.
\end{theorem}

The constant policy reproduces the no-information solution, proving
nonnegativity. A common posterior optimum proves zero value. Strict posterior
improvement integrates to positive value in the stated submodel. It does not
establish strictness when signal-contingent actions share a binding ex-ante
loss-CVaR budget. Nested action sets raise both
informed and uninformed optimized values, but their difference need not be
monotone.

\begin{proposition}[Restricted conditional complementarity]
In the posterior-separable or nonbinding-risk submodel, if actions are nested
and state--action payoff has increasing differences in
signal precision and acreage flexibility, and posterior-optimal actions differ
on events of positive probability, information and flexibility are strictly
complementary.
\end{proposition}

\begin{proposition}[Substitution through shock buffering]
Let flexibility \(\phi\) contract state-dependent margins toward a common
vector,
\[
\pi_\omega(\phi)=(1-\phi)\pi_\omega+\phi\bar\pi.
\]
At \(\phi=1\), the state no longer changes margin and information has zero
value. If information value is positive at \(\phi=0\) and the value function
is continuous, it must decrease on at least one interval; shock-buffering
flexibility substitutes for information there.
\end{proposition}

These results do not support unconditional supermodularity. For the complete
shared ex-ante loss-CVaR numerical programme we instead compute, on every adjacent
grid rectangle,
\[
\Delta_{\xi,\phi}V=
[V(\xi_2,\phi_2)-V(\xi_1,\phi_2)]
-[V(\xi_2,\phi_1)-V(\xi_1,\phi_1)].
\]
Positive, zero/boundary and negative cross-differences are numerical
classifications, not an extension of the restricted theorem. The numerical
experiments next evaluate the reversal, diversification and
information--flexibility results under a common calibration.
